% =====================================================================
% Section 5 -- Results. Steps 6.5 (5.1-5.2) and 6.6 (5.3-5.4).
% Inputs: notes 4.15-4.19 and 4.26; 4.9 points 13, 14, 16-18;
% captions.md (draft captions, to be moved into the .tex in 6.5-6.6).
% Mechanics first (JSV emphasis); engineering material condensed.
% =====================================================================
\section{Results}\label{sec:results}

\subsection{Natural frequencies and modal participation}\label{sec:frequencies}
\begin{figure}[t]
  \centering
  \includegraphics[width=\linewidth]{fig_modal}
  \caption{\TODO{fig\_modal caption (captions.md).}}
  \label{fig:modal}
\end{figure}
% Modes followed by strand; the fundamental belongs to strand B (Rayleigh)
% with bounds 4 gamma < tau_1 <= 4 tau_B; formula for F_1; intermediate drive.
\TODO{Frequencies and participation by strand; intermediate drive.}

\subsection{When the take-up mass matters}\label{sec:betaregime}
\begin{figure}[t]
  \centering
  \includegraphics[width=\linewidth]{fig_beta}
  \caption{\TODO{fig\_beta caption (captions.md).}}
  \label{fig:beta}
\end{figure}
% Generated by maps/paper_tables.py (phase 5.6) from maps/cases.py; do not edit by hand.
\begin{table*}[t]
\centering
\footnotesize
\setlength{\tabcolsep}{3.0pt}
\begin{threeparttable}
\caption{Conveyors from the literature, by decreasing length. $\xi$: take-up position from the drive exit along the return strand, in units of $L$; $\mu_r$: inertial line density of the return strand (belt and reduced idler mass); $\gamma=c_r/c_c$ at full material coupling ($\alpha=1$) and, in brackets, at the lowest coupling of the material class \citep{lodewijks2002}; $c_r$: wave speed of the return strand; $T_1$: fundamental period with the take-up mass ($\alpha=1$); $T_t$: static belt tension at the take-up; $n$: belt strands that carry the take-up pulley; $W_r=g\mu_rL$; $\beta=4M/(n^2\mu_rL)=(4/n)\lambda T_t/W_r$; $\beta_5$: take-up mass that lengthens the fundamental by 5\,\% at the case's $\xi$ and $\gamma$. Last column: what set $T_t$ according to the source.}
\label{tab:cases}
\begin{tabular}{@{}lrcrlrrrclll>{\raggedright\arraybackslash}p{24mm}@{}}
\toprule
Case & $L$ & $\xi$ & $\mu_r$ & $\gamma$ & $c_r$ & $T_1$ & $T_t$ & $n$ & $T_t/W_r$ & $\beta$ & $\beta/\beta_5$ & $T_t$ set by \\
 & (km) & & (kg/m) & & (m/s) & (s) & (kN) & & & & & \\
\midrule
Si & 13.10 & 0.001 & 35.3$^{\mathrm{c}}$ & 2.24$^{\mathrm{c}}$ (1.84) & 2171$^{\mathrm{c}}$ & 66.1$^{\mathrm{c}}$ & 115.0 & 2 & 0.025$^{\mathrm{c}}$ & 0.025$^{\mathrm{c}}$ & 0.03$^{\mathrm{c}}$ & full-load braking \\
LL & 7.60 & 0.005 & 66.9 & 1.97 (1.72) & 1649 & 46.8$^{\mathrm{a}}$ & 209.9$^{\mathrm{a}}$ & 2$^{\mathrm{a}}$ & 0.042$^{\mathrm{a}}$ & 0.084$^{\mathrm{a}}$ & 0.11$^{\mathrm{a}}$ & start grip (checked) \\
S & 7.12 & 0.001 / 0.999 & 37.8 & 1.67 (1.47) & 1659 & 40.4 / 29.0 & 173.0 & 2 & 0.066 & 0.131 & 0.22 / 0.27 & not stated (grip at $E = 2.98$) \\
H & 5.10 & 0.005 & 81.4 & 0.97$^{\mathrm{m}}$ & 1428$^{\mathrm{m}}$ & 28.4 & 49.0 & 4 & 0.012 & 0.012 & 0.03 & not stated (grip at $E \ge 4.1$) \\
G & 4.50 & 0.001 & 40.1 & 2.20 (1.77) & 1972 & 24.6$^{\mathrm{a}}$ & -- & -- & -- & 0.006$^{\mathrm{a}}$ & 0.01$^{\mathrm{a}}$ & unknown \\
SM & 3.00 & 0.100 & 72.7 & 1.95 (1.28) & 2181$^{\mathrm{c}}$ & 13.6$^{\mathrm{a}}$ & 140.0 & 2$^{\mathrm{a}}$ & 0.065 & 0.131$^{\mathrm{a}}$ & 0.17$^{\mathrm{a}}$ & not stated (model input) \\
Pa & 2.56 & 0.020 & 138.0 & 1.85 (1.30) & 2665 & 9.5 & 292.0 & 2 & 0.084 & 0.169 & 0.24 & not stated (about 2\,\% sag) \\
NC & 2.48 & 0.037 & -- & 2.46 & 1450 & 27.5 & -- & -- & -- & -- & -- & unknown \\
Lo & 1.00 & 0.001 & 21.2 & 2.76 (2.15) & 446 & 28.5 & 21.3 & 2 & 0.102 & 0.205 & 0.15 & start grip (sized) \\
Su & 0.81 & 0.010 & 39.9 & 2.58 (2.01) & 1840$^{\mathrm{c}}$ & 6.1$^{\mathrm{c}}$ & 100.0 & 2 & 0.318 & 0.635 & 0.42 & holdback \\
WR & 0.27 & -- & 31.1$^{\mathrm{c}}$ & 2.93$^{\mathrm{c}}$ (2.21) & 631$^{\mathrm{c}}$ & 5.2--5.9$^{\mathrm{a}}$ & 37.0$^{\mathrm{a}}$ & 2$^{\mathrm{a}}$ & 0.442$^{\mathrm{a}}$ & 0.884$^{\mathrm{a}}$ & 0.56--0.58$^{\mathrm{a}}$ & not stated (grip at $E = 3.0$) \\
\bottomrule
\end{tabular}
\begin{tablenotes}[flushleft]
\footnotesize
\item Superscripts: weakest provenance among the inputs of each value; c, belt or idler catalogue \citep{continental2021,fennerdunlop2009,cema}; a, assumed; m, measured. Unmarked values are published or derived from published data. The counterweight is taken as hung directly on the $n$ strands ($\lambda=1$) unless the rigging is published; for a roped rig that reduces the travel of the weight this is an upper bound of $\beta$. $E$: drive factor $e^{\mu\theta}$ implied by the published running tensions. Full provenance, notes and ranges of every input: \texttt{cases.py} in the code archive.
\item Si, \citet{sinaga2008kpc}: KPC. LL, \citet{li2009amesim}: $n$ read from the tensions of their Fig.~4. S, \citet{song2012}: take-up at the head in their Fig.~1 and at the tail in their model, where $T_t$ is published; $T_1$ and $\beta/\beta_5$ for both positions. H, \citet{harrison1983,harrison1985b}: consistency case, stepped-torque drive; $\mu_r=40$--$81$~kg/m ($\beta=0.012$--$0.024$); $\xi=0.002$--$0.02$. G, \citet{gao2026}: $\beta$ of one 1000~kg counterweight on two strands; the number of counterweights is not given and $\beta$ scales with it. SM, \citet{suchorab2025longdistance}: KGHM, variant~1; take-up between $\xi=0$ and $0.1$ ($T_1$ changes 2.6\,\%). Pa, \citet{pascual2005}: $\xi=0.005$--$0.05$ ($T_1$ changes 1.3\,\%). NC, \citet{nordell1984}: case~1; intermediate drive at 0.632$L$ from the head; no take-up data, $T_1$ with $\beta\to0$. Lo, \citet{lodewijks1996}: chapter~8. Su, \citet{surtees1995runback}: SASOL; intermediate drive at 0.197$L$ from the head, take-up after the secondary drive pulley; design~2 (design~1: $T_t=23.5$~kN, $\beta=0.149$). WR, \citet{wheatley2021}: take-up position and rigging unknown; ranges over $\xi=0.05$--$0.999$.
\end{tablenotes}
\end{threeparttable}
\end{table*}

% Threshold beta_5; identity beta = (4/n) lambda T_t/W_r = K psi tau; the
% floors that set T_t (DIN 22101 eqs. 47-52); rule T_t <~ W_r/4; claim 1
% figure ("about 1 %, at most 1.2 %...") to be confirmed in step 6.5.
\TODO{Take-up mass regime.}

\subsection{Start-up: the universal curve of a fixed--free strand}\label{sec:startup}
\begin{figure}[t]
  \centering
  \includegraphics[width=\linewidth]{fig_startup}
  \caption{\TODO{fig\_startup caption (captions.md).}}
  \label{fig:startup}
\end{figure}
\begin{figure}[t]
  \centering
  \includegraphics[width=\linewidth]{fig_crawl}
  \caption{\TODO{fig\_crawl caption (captions.md).}}
  \label{fig:crawl}
\end{figure}
% Wave law T = Z V exact for tau_a <= tau_s/2 (classic citation still to be
% found and read, 4.9 point 14); quasi-static limit; collapse of the loop;
% three formulas; fair comparison of profiles; rule tau_a >~ 2-3 tau_1;
% initial crawl (notes 4.17).
\TODO{Universal start-up curve.}

\subsection{Validity: slack, take-up and drive grip}\label{sec:validity}
\begin{figure}[t]
  \centering
  \includegraphics[width=\linewidth]{fig_validity}
  \caption{\TODO{fig\_validity caption (captions.md).}}
  \label{fig:validity}
\end{figure}
% Three conditions (notes 4.18); grip condition with the start factor from
% the universal curve; DIN 22101 sec. 7.3.1 and 8.2.2 (4.9 point 16).
\TODO{Validity conditions.}
